\section{The semantic core: game forms and kernel games}
\label{sec:semantic-core}

The library is organized around two records — \TT{GameForm} and
\TT{KernelGame} — that together act as the universal target into
which every other representation compiles, and the universal
domain on which every solution concept is defined.  This section
introduces them, the morphisms between them, and the
deviation-family schema that unifies the equilibrium notions of
\S\ref{sec:solution-concepts}.

\subsection{Game forms}
\label{sec:semantic-core:gameform}

A \emph{game form} over a player type $\Players$ is a triple
\begin{equation}
\label{eq:gameform}
  \GF \;=\; \bigl(\,\Strat : \Players \to \mathsf{Type},\;
                    \Outcome : \mathsf{Type},\;
                    \K : {\textstyle\prod_i}\Strat_i \to \PMF{\Outcome}\,\bigr).
\end{equation}
A \emph{strategy profile} is an element of $\prod_i \Strat_i$,
written $\Profile$.  The kernel $\K$ assigns to each profile a
distribution over outcomes; outcome distributions and correlated
outcome distributions are derived in the obvious way:
\[
  \mathrm{outcomeDist}(\GF, \Profile) \;=\; \K(\Profile),
  \qquad
  \mathrm{correlatedOutcome}(\GF, \mu) \;=\; \mu \mathbin{\bind} \K
\]
where $\mu \in \PMF{\prod_i \Strat_i}$ is a correlated profile
distribution and $\bind$ is the \TT{PMF} bind.  For a finite
$\Players$ with \TT{DecidableEq}, unilateral deviation is the
ordinary update operation
\[
  \Profile[i \mapsto s'] \;\coloneqq\;
    \Profile \text{ with coordinate } i \text{ replaced by } s'.
\]
Two derived families of distributions formalize the kinds of
deviation that equilibrium notions quantify over.  The
\emph{constant} deviation family
\[
  \mathrm{constDev}(\GF, \mu, i, s') \;=\;
    \mu \mathbin{\bind} \bigl(\Profile \mapsto \Dirac_{\,\Profile[i \mapsto s']}\bigr)
\]
replaces $i$'s coordinate by the same $s'$ regardless of what $\mu$
sampled.  The \emph{recommendation-dependent} family
\[
  \mathrm{depDev}(\GF, \mu, i, \mathrm{dev}) \;=\;
    \mu \mathbin{\bind} \bigl(\Profile \mapsto
      \Dirac_{\,\Profile[i \mapsto \mathrm{dev}(\Profile_i)]}\bigr)
\]
lets the deviation be a function of $i$'s recommendation; this is
the family used in correlated equilibrium~\citep{Aumann1974}.

\designnote{The two deviation families are kept as named records
rather than inlined into each equilibrium notion.  This is the
single point at which the difference between Nash, coarse
correlated equilibrium, and correlated equilibrium becomes
visible:\ Nash and CCE quantify over $\mathrm{constDev}$, CE
quantifies over $\mathrm{depDev}$.  We exploit this in
\S\ref{sec:solution-concepts:family} by stating monotonicity
results once on the schema.}

\leancorr{\TT{Core/GameForm.lean} for the record and derived
operations; supplement~§2.1 has the full Lean signatures.}

\subsection{Kernel games}
\label{sec:semantic-core:kernelgame}

A \emph{kernel game} adds a utility to a game form,
\begin{equation}
\label{eq:kernelgame}
  \KG \;=\; \bigl(\,\Strat,\;\Outcome,\;\util : \Outcome \to \Payoff{\Players},\;\K\,\bigr),
\end{equation}
where $\Payoff{\Players} \coloneqq \Players \to \RR$.  The
\emph{expected utility} of player $i$ under profile $\Profile$ is
\begin{equation}
\label{eq:eu}
  \EU(\KG, \Profile, i) \;=\; \sum_{\omega} \K(\Profile)(\omega)\;\util(\omega)(i),
\end{equation}
a finite sum under the finite-outcome hypotheses used by the major
theorems.  At the core API level the implementation is mathlib's
unordered-sum expectation for $\PMF{\cdot}$, well-defined on
bounded utilities over an arbitrary (countable) \Outcome; the
finite-support lemmas recover the displayed formula, and a
parallel family of \TT{\_of\_bounded} lemmas replays the same
algebra without a \TT{[Finite \Outcome]} hypothesis.  Expected utility is the only place
where utility enters a numerical inequality.

The conversion $\toGF : \KG \mapsto \GF$ forgets utility; the
inverse $\toKG_u : \GF \mapsto \KG$ adjoins a utility $u$.  These
are mutual inverses up to record equality.

A kernel game admits an explicit \emph{mixed extension}:\ each
player's strategy type is replaced by $\PMF{\Strat_i}$, and the
outcome kernel samples a pure profile from the independent product
$\bigotimes_i \Profile_i \in \PMF{\prod_i \Strat_i}$ and applies
the original kernel.  Only the player index $\Players$ needs to
be finite (so the product over players is well-defined); the
per-player strategy types $\Strat_i$ and the outcome type
$\Outcome$ may be countably infinite.  The mixed extension is
itself a kernel game, indexed by the same $\Players$, and serves
as the carrier of mixed strategies in Nash existence
(\S\ref{sec:theorems:nash-mixed}), where the existence proof
adds finiteness back as a hypothesis on $\Strat_i$ and $\Outcome$.

\leancorr{\TT{Core/KernelGame.lean} (\TT{KernelGame},
\TT{KernelGame.eu}, \TT{KernelGame.mixedExtension}, \TT{reindex});
\TT{Core/GameForm.lean} (\TT{toGameForm}, \TT{withUtility} and the
round-trip lemmas); \TT{Math/PMFProduct.lean} (\TT{pmfPi});
supplement~§2.2--2.4.}

\subsection{Solution concepts as predicates: the \texttt{*For} style}
\label{sec:semantic-core:for-style}

The library follows a uniform style for solution concepts that
rewards a brief explanation, since it is unusual.  Each concept
exists in two forms:
\begin{itemize}
\item \emph{EU-specific.} $\IsNash$, $\IsBR$,
$\WDom$, $\SDom$, $\IsCorrelatedEq$, $\IsCoarseCorrelatedEq$:
  predicates on $\KG$ phrased directly in terms of the expected
  utility \eqref{eq:eu}.

\item \emph{Preference-parameterized.} $\IsNashFor$,
  $\IsBRfor$, $\WDomFor$, $\SDomFor$,
  $\IsCorrelatedEq_{\mathit{For}}$, $\ldots$:\ predicates on $\GF$
  taking an additional argument, a preference relation
  $\Pref : \Players \to \PMF{\Outcome} \to \PMF{\Outcome} \to \mathsf{Prop}$
  on outcome distributions.  ``$\Pref_i \,d_1\,d_2$'' means
  ``player $i$ weakly prefers distribution $d_1$ over $d_2$.''
\end{itemize}
The two are tied together by lemmas of the shape
\[
  \KG.\IsNash(\Profile) \;\Longleftrightarrow\;
    \KG.\toGF.\IsNashFor\bigl(\Pref^{\EU}_{\KG},\,\Profile\bigr),
\]
where $\Pref^{\EU}_{\KG,\,i}\,d_1\,d_2 \coloneqq \mathbb{E}_{d_1}[\util(\cdot)(i)] \ge
\mathbb{E}_{d_2}[\util(\cdot)(i)]$, the expected-utility weak
preference induced by $\KG$.

\paragraph{Why two layers?}  Two reasons.  First, several
classical results — Kuhn's theorem, the existence of correlated
equilibria, the dominance$\,\Rightarrow\,$Nash chain — are
\emph{purely structural}: they hold whenever players have a
preorder on outcome distributions, regardless of whether that
preorder comes from expected utility, lexicographic preferences,
or maximin attitudes.  Stating them at the $\GF.\IsNashFor$ layer
keeps them generic and stops them from masquerading as
EU-specific theorems.  Second, the EU-specific layer keeps the
common case ergonomic:\ users who care only about expected
utility never have to mention preference relations.

\paragraph{Why explicit-argument preferences and not a typeclass?}
This is a question worth posing directly, since at first glance a
typeclass \TT{[Preference G]} carrying the preference relation —
with a default \TT{euPref} instance — would make the EU-only user's
syntax uniform (\TT{G.IsNash}~\Profile\ in every case).  We considered
this and rejected it because a substantial fraction of the
library's theorems quantify over \emph{more than one} preference
simultaneously, and typeclasses fit such theorems poorly:

\begin{itemize}\itemsep2pt
\item Monotonicity (\TT{IsNashFor.mono}, \TT{IsDominantFor.mono})
  takes two preferences $\Pref^{(1)}, \Pref^{(2)}$ together with a
  pointwise-implication hypothesis between them, and concludes a
  predicate for $\Pref^{(1)}$ from a predicate for $\Pref^{(2)}$.

\item Pareto dominance is parameterized by a \emph{pair} $(\Pref,
  \sPref)$ of weak-and-strict preferences satisfying compatibility
  conditions; \TT{ParetoDominatesFor pref spref σ τ} takes both
  explicitly.

\item The transport theorems of
  \S\ref{sec:semantic-core:transport} relate preferences across
  two presentations of a game, either by pullback along an outcome
  bijection or by comparison after both outcome laws are projected
  to a common observation space.  In both cases two preferences
  appear in the same statement, one a function of the other.

\item Strict-from-weak / weak-from-strict bridges of
  \S\ref{sec:solution-concepts:dominance} relate a strict preference
  to a weak preference under a hypothesis $\sPref \Rightarrow \Pref$.
\end{itemize}

Each of these binds two preferences at once.  A typeclass-based
design would either force users to construct ad-hoc local instances
for every multi-preference argument and switch between them by
\TT{let} bindings (noisy and error-prone), or fall back to
explicit arguments for these theorems anyway (defeating the
uniformity the typeclass was supposed to deliver).  Explicit
arguments are the natural shape from the start, and we add the
EU-specific predicates as a parallel surface so that the EU-only
ergonomic is recovered without invoking typeclass inference at
all.

There is one place where the library does use a typeclass on
preferences:\ \TT{PrefPreorder pref} asserts that a preference
relation is a preorder (reflexive and transitive).  This is the
correct use of a typeclass — it expresses a \emph{property} of a
given relation, not a choice between relations, and Lean's
inference resolves it without any of the diamond/multi-preference
issues above.

\tradeoff{The cost of the two-layer presentation is one extra
indirection for non-EU users (they must instantiate the preference
explicitly) and a handful of \texttt{\_iff\_} bridges to keep the
EU-specific definitions in step with their preference-parameterized
counterparts.}

\subsection{Morphisms, simulations, and isomorphisms}
\label{sec:semantic-core:morphisms}

Game forms and kernel games each support a notion of structured
map closed under composition with identities, with the standard
identity- and associativity laws proved in the library.  We treat
them as ordinary categories of presentation but do not exploit
further categorical structure in the chapter.  At the game-form
level, a \emph{protocol map} from $\GF \to \GF'$ is a triple
\[
  (\stratMap, \outcomeMap, \mathrm{preserved})
  \quad\text{with}\quad
  \stratMap : \forall i, \Strat_i \to \Strat'_i,\;
  \outcomeMap : \Outcome \to \Outcome',
\]
satisfying the \emph{outcome-kernel preservation} law
\begin{equation}
\label{eq:protocolmap}
  \K'(\stratMap \circ \Profile) \;=\;
    \K(\Profile) \mathbin{\bind} \bigl(\omega \mapsto \Dirac_{\,\outcomeMap(\omega)}\bigr).
\end{equation}
This is the discrete-Giry analogue of a coalgebra homomorphism,
and is what makes a bridge between two languages strategically
faithful (i.e.\ Nash equilibria pull back, dominance pulls back,
etc., under hypotheses on $\stratMap$).  At the \TT{KernelGame}
level we add the utility-preservation condition
$\util'(\outcomeMap(\omega))(i) = \util(\omega)(i)$, yielding a
\emph{kernel-game simulation}, and call the morphism an
\emph{EU-isomorphism} when both $\stratMap$ and $\outcomeMap$ are
bijective and the inverse is also a simulation.

\leancorr{Morphism records in \TT{Core/GameMorphism.lean}
(\TT{ProtocolMap}, \TT{KernelGame.Morphism}, with the
\TT{id\_comp}/\TT{comp\_id}/\TT{comp\_assoc} laws);
\TT{Core/GameSimulation.lean} for kernel-game simulations and
bisimulations; \TT{Core/GameIsomorphism.lean} for isomorphisms;
\TT{Core/UtilityInvariance.lean} packages the standard transport
results; \TT{Concepts/EUProperties.lean} and
\TT{Concepts/OfEUProperties.lean} record EU structural properties
and the \TT{ofEU} round-trip.  Supplement~§2.6.}

\subsection{Reindexing and the player type}
\label{sec:semantic-core:reindex}

The player type $\Players$ is left abstract throughout, but
bridges between languages typically pick a canonical finite
player type — most often $\mathrm{Fin}\,n$ for the cardinality of
the source's player set.  The library provides
$\mathrm{reindex} : (\Players \simeq \kappa) \to \KG_\kappa \to
\KG_\Players$, transporting strategy spaces, utilities, and the
outcome kernel along an equivalence.  This is presentation
bookkeeping, but it is the kind of bookkeeping that, omitted, leaks
$\mathrm{Fin}\,n$ assumptions into every downstream consumer.  The
chapter writes $\KG[e]$ informally for the reindexed game.

\subsection{Transport of solution concepts under bridges}
\label{sec:semantic-core:transport}

The use we make of protocol maps and EU-isomorphisms in
\S\ref{sec:languages:bridges} relies on transport theorems that
turn semantic preservation into equilibrium preservation.  There
are two levels.  The first is the familiar map-based case: a
bijective relabeling of strategies and outcomes preserves every
deviation menu exactly.  The second is more general and is the one
needed for representation changes such as richer traces, serialized
moves, or Kuhn-style realization relations: profiles need only have
the same \emph{observed} outcome law, and target-side deviations
must be simulated by source-side deviations with the same observed
law.

\begin{theorem}[Solution-concept transport]
\label{thm:transport}
Let $\GF, \GF'$ be game forms over the same player type
$\Players$ and let $f : \GF \to \GF'$ be a protocol map with
$\stratMap_i$ bijective for every $i$ and $\outcomeMap$ bijective.
Let $\Pref'$ be a preference on $\PMF{\Outcome'}$ and let
$\Pref_i\,d_1\,d_2 \coloneqq \Pref'_i\,(\outcomeMap_*\,d_1)\,(\outcomeMap_*\,d_2)$
denote its pullback along the outcome bijection.  Then for every
profile $\Profile$,
\[
  \GF.\IsNashFor(\Pref, \Profile)
  \;\Longleftrightarrow\;
  \GF'.\IsNashFor\bigl(\Pref',\,\stratMap \circ \Profile\bigr).
\]
The analogous statements hold for $\IsBRfor$, $\IsDomFor$,
$\WDomFor$, $\SDomFor$, $\IsCorrelatedEqFor$, and
$\IsCoarseCorrelatedEqFor$.  Specializing $\Pref$ to expected
utility under a kernel-game simulation that preserves utility
yields the EU-specific transport laws.
\end{theorem}

\begin{theorem}[Observed-law deviation transport]
\label{thm:observed-deviation-transport}
Let $\GF$ and $\GF'$ be game forms over the same player type, let
$\Omega$ be a common observation space, and let
$V : \Outcome \to \Omega$ and $V' : \Outcome' \to \Omega$ be
outcome views.  Suppose a relation $R(\Profile,\Profile')$ between
profiles satisfies
\[
  V_*\,\K(\Profile) = V'_*\,\K'(\Profile')
\]
whenever $R(\Profile,\Profile')$ holds.  Suppose further that for
every target unilateral deviation $\Profile'[i \mapsto s'_i]$
there exists a source unilateral deviation
$\Profile[i \mapsto s_i]$ with the same observed outcome law:
\[
  V_*\,\K(\Profile[i \mapsto s_i])
  =
  V'_*\,\K'(\Profile'[i \mapsto s'_i]).
\]
Then, for every preference $\Pref^\Omega$ on $\PMF{\Omega}$,
\[
  \GF.\IsNashFor(V^*\Pref^\Omega,\Profile)
  \;\Longrightarrow\;
  \GF'.\IsNashFor((V')^*\Pref^\Omega,\Profile').
\]
The same statement holds at the level of arbitrary
\TT{ProfileDeviationFamily} equilibria: if every target-family
deviation is matched by a source-family deviation with the same
observed law, then $\mathcal{D}$-equilibrium transfers from the
source presentation to the target presentation.  The Nash theorem is
the constant-deviation, point-mass specialization.
\end{theorem}

\noindent
The direction of the deviation-simulation hypothesis is important.
To transfer Nash from $\GF$ to $\GF'$, one must rule out every
possible deviation in $\GF'$; therefore each such target deviation
must be represented by some source deviation whose observed law the
source Nash hypothesis already controls.  Equality of the base
outcome law alone is not enough, because Nash quantifies over
counterfactual unilateral deviations.

\noindent
This is the operational content of contributions 3 and 5 in
\S\ref{sec:introduction}:\ a bridge between two languages, viewed
as a protocol-preserving (resp.\ EU-preserving) map between the
compiled game forms (resp.\ kernel games), automatically
transports solution concepts.  Concrete bridges
(\S\ref{sec:languages:bridges}) instantiate the theorem with
specific maps and inherit a Nash/dominance/correlated-eq
preservation law for free.

\leancorr{\TT{Core/UtilityInvariance.lean} and
\TT{Core/GameSimulation.lean} for the map-based transport layer;
\TT{Concepts/DeviationSimulation.lean} for
\TT{OutcomeView}, \TT{NashDeviationSimulation}, and
\TT{DeviationFamilySimulation}.}

\subsection{Designated facts}
\label{sec:semantic-core:designated}

Three theorems on $\GF$ deserve early visibility:
\begin{itemize}\itemsep2pt
\item \emph{Dominance implies Nash}:\ for any preference
$\Pref$, if every player's strategy in $\Profile$ is
$\Pref$-dominant, $\Profile$ is a $\Pref$-Nash equilibrium.

\item \emph{Best-response characterization}:\ $\Profile$ is
$\Pref$-Nash iff each $\Profile_i$ is a $\Pref$-best response
against the rest.

\item \emph{Correlated $\Rightarrow$ coarse correlated}:\ every
correlated equilibrium is a coarse correlated equilibrium.
\end{itemize}
None of these depends on EU; they hold for any preference
preorder.  The EU-specific versions are corollaries of the
generic statements via the EU/preference iff-bridges of
\S\ref{sec:semantic-core:for-style}.

\leancorr{All three in \TT{Core/GameForm.lean}
(\TT{dominant\_is\_nash\_for},
\TT{isNashFor\_iff\_bestResponseFor},
\TT{IsCorrelatedEqFor.toCoarseCorrelatedEqFor}).}
